% !PRINTED 162
\weberpage{162}
%<<CONT
eine beliebige Form $\varphi(x)$ der $n$ Veränderlichen $x$ zu nehmen,
die Functionen:
\weq{2}{\varphi[A(x)],\ \varphi[B(x)],\ \varphi[C(x)],\ \ldots}
für alle Substitutionen der Gruppe zu bilden, und irgend eine
symmetrische Function der Formen (2) für $\Phi$ zu nehmen.

Denn wendet man auf $(x)$ eine Substitution der Gruppe $S$
an, so ändert sich die Gesammtheit der Functionen (2) nicht; es
wird nur ihre Reihenfolge eine andere.

Man kann den Begriff der Invarianten noch allgemeiner
fassen, wie folgt:

2. \gesperrt{Eine Form $\mathit\Psi(x_1,x_2,\ldots,x_n)$ heisst auch dann eine
Invariante der Gruppe $S$, wenn sie constante
Factoren annimmt, wenn auf die Variablen $(x)$
die Substitutionen $A,B,C \ldots$ der Gruppe $S$
angewandt werden.}

Durch Formeln wird diese Eigenschaft so ausgedrückt:
\weq{3}{\mathit\Psi[A(x)]=\alpha\,\mathit\Psi(x),\ \mathit\Psi[B(x)]=\beta\,\mathit\Psi(x),\ \mathit\Psi[C(x)]=\gamma\,\mathit\Psi(x)\ldots}
worin die Co\"efficienten $\alpha,\beta,\gamma\ldots$ von den $x$ unabhängig sind.

Wenn eine Unterscheidung n\"othig ist, wollen wir die in 1.
definirten Formen \gesperrt{absolute Invarianten}, und die in 2. \gesperrt{relative
Invarianten} der Gruppe $S$ nennen.

Aus den Formeln (3) ergiebt sich:
\weq{4}{\mathit\Psi[AB(x)] = \alpha\beta\,\mathit\Psi(x),}
und daraus folgt, dass die Factoren $\alpha,\beta,\gamma,\ldots$ in ihrer Gesammtheit,
bei der Zusammensetzung durch Multiplication, eine Gruppe
bilden müssen.

Zwischen dieser (commutativen) Gruppe und der Gruppe $S$
besteht ein im Allgemeinen mehrstufiger Isomorphismus, da
verschiedene Substitutionen aus $S$ zu demselben Factor $\alpha$ führen
können.

Ist $\mu$ der Grad der Gruppe $S$, so ist der Grad eines jeden
Elementes $A, B \ldots$ von $S$ ein Theiler von $\mu$, und folglich ist
$A^\mu = B^\mu \cdots$ gleich der identischen Substitution.

Hieraus folgt, dass die Factoren $\alpha,\beta,\gamma\ldots$ \gesperrt{$\mu^{\mathrm{te}}$ Einheitswurzeln} sind.

Ist $e$ die kleinste positive, der Bedingung
\weq{5}{\alpha^e = \beta^e = \gamma^e = \cdots = 1}
genügende Zahl, so sind die $\alpha, \beta, \gamma, \ldots$ zugleich $e^{\mathrm{te}}$ Einheits
%CONT>>
